\documentclass{article}
\usepackage[margin=25mm]{geometry}
\usepackage{amsmath,amssymb}
\usepackage{xurl}
\usepackage[colorlinks=true,linkcolor=blue,citecolor=blue,urlcolor=blue]{hyperref}
% Shared commands for notes in docs/paper-gaps/.

\DeclareUnicodeCharacter{2080}{\ifmmode{}_{0}\else\textsubscript{0}\fi}
\DeclareUnicodeCharacter{2081}{\ifmmode{}_{1}\else\textsubscript{1}\fi}
\DeclareUnicodeCharacter{2082}{\ifmmode{}_{2}\else\textsubscript{2}\fi}
\DeclareUnicodeCharacter{2099}{\ifmmode{}_{n}\else\textsubscript{n}\fi}

\newcommand{\Lean}{\textsc{Lean}}
\ExplSyntaxOn
\NewDocumentCommand{\leanid}{m}{
  \ifmmode
    \textsf{\detokenize{#1}}
  \else
    \group_begin:
    \urlstyle{sf}
    \tl_set:Nn \l_tmpa_tl {#1}
    \tl_replace_all:Nnn \l_tmpa_tl {\_} {_}
    \exp_args:NV \path \l_tmpa_tl
    \group_end:
  \fi
}
\ExplSyntaxOff
\providecommand{\tr}{\operatorname{tr}}
\newcommand{\tnleanIssueBase}{https://github.com/LionSR/TNLean/issues/}
\newcommand{\tnleanPrBase}{https://github.com/LionSR/TNLean/pull/}
\newcommand{\tnleanDocBase}{https://sirui-lu.com/TNLean/docs/}
\newcommand{\ghissue}[1]{\href{\tnleanIssueBase#1}{GitHub issue \##1}}
\newcommand{\ghpr}[1]{\href{\tnleanPrBase#1}{PR \##1}}
\newcommand{\doclink}[1]{\href{\tnleanDocBase#1.html}{\path{#1}}}

% Dirac notation and the identity operator, as in blueprint/src/macros/common.tex.
% \providecommand keeps note-local definitions authoritative.
\usepackage{dsfont}
\providecommand{\ket}[1]{|#1\rangle}
\providecommand{\bra}[1]{\langle #1|}
\providecommand{\braket}[2]{\langle #1 | #2 \rangle}
\providecommand{\ketbra}[2]{|#1\rangle\!\langle #2|}
\providecommand{\Id}{\mathds{1}}
\providecommand{\one}{\mathds{1}}
\providecommand{\C}{\mathbb{C}}
\providecommand{\MN}[1]{M_{#1}(\C)}
\providecommand{\MD}{M_D(\C)}

% Verdict marker: \gapnote{<kind>}{<status>}. Kind is the policy
% classification (clarification, local-correction, scope-restriction,
% unfaithful, false-source, open-gap); status is open, wip (elimination
% underway), resolved, or historical. Machine-read by the site index;
% prints a verdict line.
\newcommand{\gapnote}[2]{\par\noindent\textbf{Verdict:} #1 (#2).\par}

\providecommand{\gapnote}[2]{\par\noindent\textbf{Verdict:} #1 (#2).\par}
\title{Support Size and Negative Local Terms in the ALVV Cluster Approximation}
\author{}
\date{October 25, 2026}
\begin{document}
\maketitle
\gapnote{false-source}{open}

\section{The printed estimate}

In the section on soft truncation, Arad, Landau, Vazirani and
Vidick~\cite{Arad2016RigorousRG},
\href{https://arxiv.org/abs/1602.08828v2}{arXiv:1602.08828v2}, state
Theorem~\path{thm:cluster-approx} for
$H=\sum_{i=1}^{n-1}h_i$ on a chain of $n$ qudits, with $h_i$ supported on
$\{i,i+1\}$ and $\|h_i\|\le1$. The paragraph defining the expansion explicitly
allows terms that are not nonnegative. For a word $w$ in the interaction
labels, put
\begin{align}
    f(w)=\frac{(-\beta)^{|w|}}{|w|!}h_w,
    \qquad h_w=h_{w_1}\cdots h_{w_{|w|}}.
    \label{eq:alvv_cluster_word}
\end{align}
The operator $M_r(H)$ keeps the words whose support has connected components
of fewer than $r$ physical sites. The printed assertion is
\begin{align}
    \|e^{-\beta H}-M_r(H)\|
    \le \exp\bigl(n^2(e^\beta-1)^r\bigr)-1,
    \qquad 0<\beta<\log2.
    \label{eq:alvv_cluster_printed}
\end{align}
The source locations are \path{local_alg_arxiv_v2.tex}, lines 1270--1291
(definition and statement), and the appendix \emph{Constructing an MPO for the
    cluster expansion}, lines 1924--1968 (proof).
There are two independent obstructions: the proof counts interaction labels
where the statement counts physical sites, and it bounds a complementary
exponential by one without assuming positivity.

\section{Two sites: the support cutoff}

Take two qubits and the single term
$H=h_1=P=|00\rangle\langle00|$, with
$r=2$ and $\beta=\log(11/10)$. This term is positive, has norm one,
and acts nontrivially on both qubits. Every nonempty word evaluates to $P$,
so its connected support has two sites and it is excluded. The empty word
is retained. Thus
\begin{align}
    M_2(H)                  & =I,                 &
    e^{-\beta P}            & =I+(e^{-\beta}-1)P, &
    \|e^{-\beta P}-M_2(H)\| & =\frac1{11}.
    \label{eq:alvv_cluster_two_sites}
\end{align}
The right side of (\ref{eq:alvv_cluster_printed}) is
$\exp(4(1/10)^2)-1=\exp(1/25)-1$. For $0<x<1$, comparison of the
exponential and geometric series gives $e^x<1/(1-x)$. Consequently
\begin{align}
    \exp(1/25)-1<\frac1{24}<\frac1{11}.
    \label{eq:alvv_cluster_two_sites_comparison}
\end{align}
This contradicts the printed bound even for positive local terms. A connected
interval of $r$ sites requires only $r-1$ distinct nearest-neighbour
interaction labels. The factor $(e^\beta-1)^{|I|}$ in the appendix counts
labels, not sites.

\section{Forty sites: negative local terms}

Changing the cutoff to count interaction labels does not resolve the
unrestricted statement. On forty qubits let
\begin{align}
    h_i   & =-P_i,                                                          &
    P_i   & =|00\rangle\langle00|_{i,i+1}\otimes I_{\{i,i+1\}^{\mathrm c}},
    \qquad 1\le i\le39,
    \label{eq:alvv_cluster_negative_terms}                                             \\
    \beta & =\log(3/2),                                                     & r & =20.
    \notag
\end{align}
These are commuting Hermitian terms of norm one with two-site
supports. Let $\Omega=|0\rangle^{\otimes40}$. For every word of length $k$,
\begin{align}
    f(w)=\frac{\beta^k}{k!}P_{w_1}\cdots P_{w_k},
    \qquad \langle\Omega,f(w)\Omega\rangle=\frac{\beta^k}{k!}\ge0.
    \label{eq:alvv_cluster_positive_coefficients}
\end{align}
Every nonempty product of the $P_i$ fixes precisely the qubits in the union
of its edges to zero and is the identity elsewhere. Its actual physical
support is therefore that union; there is no cancellation of supports.

Consider only words containing each of labels $1,\ldots,20$ at least once,
with labels $21,\ldots,39$ unrestricted. Each such word has a connected
component with at least twenty interaction labels and twenty-one physical
sites. It is omitted for either interpretation of the cutoff $r=20$.
Moreover,
\begin{align}
    \sum_w\|f(w)\|
    \le\sum_{k=0}^{\infty}\frac{39^k\beta^k}{k!}=e^{39\beta},
    \label{eq:alvv_cluster_absolute_convergence}
\end{align}
so selecting and regrouping these words is justified. At fixed
multiplicities $m_1,\ldots,m_{39}$, their number is
$(\sum_i m_i)!/\prod_i m_i!$. Hence their total expectation is
\begin{align}
    \prod_{i=1}^{20}\Bigl(\sum_{m_i\ge1}\frac{\beta^{m_i}}{m_i!}\Bigr)
    \prod_{i=21}^{39}\Bigl(\sum_{m_i\ge0}\frac{\beta^{m_i}}{m_i!}\Bigr)
     & =(e^\beta-1)^{20}e^{19\beta}
    =\frac{3^{19}}{2^{39}}.
    \label{eq:alvv_cluster_selected_weight}
\end{align}
All omitted words have nonnegative expectation by
(\ref{eq:alvv_cluster_positive_coefficients}). Since $\|\Omega\|=1$,
\begin{align}
    \|e^{-\beta H}-M_{20}(H)\|
    \ge\langle\Omega,(e^{-\beta H}-M_{20}(H))\Omega\rangle
    \ge\frac{3^{19}}{2^{39}}.
    \label{eq:alvv_cluster_negative_lower_bound}
\end{align}
This also holds for the alternative truncation by interaction-label
components. Put $x=40^2(1/2)^{20}=25/16384$. The proposed upper bound satisfies
\begin{align}
    e^x-1<\frac{x}{1-x}=\frac{25}{16359}<\frac{3^{19}}{2^{39}},
    \label{eq:alvv_cluster_negative_comparison}
\end{align}
where the final inequality is the exact integer comparison
\begin{align}
    25\cdot2^{39}=13743895347200
    <19013435338653=16359\cdot3^{19}.
    \label{eq:alvv_cluster_integer_comparison}
\end{align}
Equations (\ref{eq:alvv_cluster_negative_lower_bound}) and
(\ref{eq:alvv_cluster_negative_comparison}) contradict
(\ref{eq:alvv_cluster_printed}) under either cutoff convention.

\newpage
\section{The missing hypotheses and the remaining problem}

The appendix uses $\|e^{-\beta H_{\overline I}}\|\le1$ at line 1948.
This follows for positive local terms, but not from Hermiticity and
$\|h_i\|\le1$. If the complementary sum has $m$ terms, these hypotheses
give only $\|e^{-\beta H_{\overline I}}\|\le e^{\beta m}$. The general
positive-term convention earlier in the paper (lines 267, 806 and 1175)
cannot supply the missing assumption here: lines 1271--1272 explicitly
allow terms that are not nonnegative.

A corrected theorem must specify its cutoff and its energy dependence.
For positive terms with an interaction-label cutoff, the occupancy factor
$(e^\beta-1)^r$ has the intended meaning. For a physical-site cutoff, an
interval of $r$ sites gives only $r-1$ mandatory labels. For arbitrary
Hermitian terms an additional bound on the complementary exponential is
needed. Each alternative requires its own statement and proof. Adding
positivity alone, or changing the cutoff alone, does not repair
(\ref{eq:alvv_cluster_printed}). Shifting the Hamiltonian by a scalar also
requires checking the induced change of the word-defined truncation.

The printed theorem is not formalized in TNLean. The existing Hastings
cluster estimates retain the partition function and the weight
$(e^{\beta J}-1)e^{SJ\beta}$; they do not assert
(\ref{eq:alvv_cluster_printed}). Their formalized statements remain
unaffected. The efficient soft-truncation approximation in
Lemma~\texttt{lem:exp-approx-eff}, whose proof invokes this theorem, needs
a corrected analytic estimate before that argument can be formalized.
The exact matrix product operator construction of the selected expansion
is a separate algebraic question. The correction is tracked by
\href{https://github.com/LionSR/TNLean/issues/9854}{issue 9854} and
\href{https://github.com/LionSR/TNLean/issues/9847}{issue 9847}.

\bibliographystyle{unsrt}
\bibliography{references}
\end{document}
